\documentclass[11pt]{article}
\usepackage[margin=1in]{geometry}
\usepackage{amsmath,amssymb,amsthm,hyperref}
\newtheorem{theorem}{Theorem}
\newtheorem{proposition}[theorem]{Proposition}
\newtheorem{definition}[theorem]{Definition}
\newcommand{\NN}{\mathbb N}
\newcommand{\RR}{\mathbb R}
\newcommand{\Q}{\mathbb Q}
\newcommand{\dd}{\mathrm d}
\title{Future-safe quotients and source-sensitive obstructions\\\large An Observational SUCC research draft}
\author{Observational SUCC research program}
\date{18 October 2026}
\begin{document}
\maketitle
\begin{abstract}
We give a minimal deterministic contract for discarding information while
preserving declared future reports, partial actions and costs. Predictive
sufficiency is separated from autonomous quotient dynamics, source-realizable
paths from compatible report paths, and existential lifts from online
realization. The quotient and finite path mechanisms specialize classical
automata and coalgebra theory. A self-contained arithmetic construction shows
that any finite ordinary-prime prefix, density one, a prime number theorem and
the critical-strip zero divisor together fail to identify successor--dilation
rigidity within a generalized-prime source class. Finally, we derive a
boundary-aware resolvent estimate for the actual real-Gamma tail used in the
signed Suzuki form. Six scoped quotient/path lemmas have been kernel checked;
the arithmetic and analytic proofs are paper proofs with exact finite controls.
Academic priority is unverified. No first-window Weil sign or RH proof is claimed.
\end{abstract}

\section{Minimal specification and proof authority}
A deterministic contract is a set $X$ of complete source/history states, an
action set $A$, a report map $o:X\to O$, and partial maps
$t_a:X\rightharpoonup X\times\RR_{\ge0}$. Reports may include current
source tags and the outputs of a declared proof verifier. An immutable
provenance label, when required, is separately preserved by every action.
Epochs and current source configurations can change inside $X$.
A state need not admit a finite effective representation. A hypothetical
path is not an actual event journal, and equality of verifier outputs does
not establish soundness of that verifier.

A finite structural trace records its initial report, all reached reports,
costs and the first unavailable action. Unavailability becomes an observed
separator only under a specified lawful refusal/menu instrument. The
statements below concern this declared structural contract; they do not
authenticate a physical source or provide an undeclared instrument.

\section{Exactly what can be forgotten}
Write $x\simeq y$ when all finite structural traces agree. For a representation
$q:X\twoheadrightarrow Y$, call $q$ \emph{predictively sufficient} if every
finite trace is a function of $q(x)$. Call it \emph{autonomous} if reports and
partial cost-labelled actions descend to $Y$.

\begin{theorem}[Classical quotient, with explicit contract]
Predictive sufficiency is equivalent to $\ker q\subseteq\simeq$.
Autonomy is equivalent to the following local laws whenever $q(x)=q(y)$:
\begin{enumerate}
\item $o(x)=o(y)$;
\item $t_a(x)$ is defined exactly when $t_a(y)$ is defined;
\item if $t_a(x)=(x',c)$ and $t_a(y)=(y',d)$, then $c=d$ and $q(x')=q(y')$.
\end{enumerate}
The relation $\simeq$ is the greatest equivalence satisfying these laws.
Its quotient is autonomous. Every surjective autonomous representation
$q$ admits a unique map $f:Y\to X/{\simeq}$ such that $\pi=fq$.
Every executable quotient word lifts from every source representative,
preserving each charged cost and the final quotient state.
\end{theorem}
\begin{proof}
A function factors through a surjection exactly when it is constant on its
fibers; apply this to each trace. If actions descend, the three local laws
follow by comparing their images. Conversely define the quotient output
and action using a representative; the laws make the definition independent
of that choice. Comparing empty traces gives output equality. Comparing
one-letter prefixes and then arbitrary suffixes gives availability, cost
and successor equivalence for $\simeq$. Conversely the local laws propagate
trace equality by induction on word length. Thus every such equivalence is
contained in $\simeq$. Set $f(q(x))=[x]$; this is well-defined, unique and
commutes with the induced actions. Finally, a quotient transition must equal
the image of the transition from any representative, so it has a lift from
that representative; iterate over the word.
\end{proof}

This is a specialization of final semantics and quotient descent for
$F(X)=O\times(1+X\times\RR_{\ge0})^A$ \cite{rutten,silva}; histories and
source labels do not by themselves establish novelty. The repository's B15
and causal Nerode development already establish the total deterministic core.
The constructive Lean theorem takes a supplied section of $q$ and does not
invoke choice. Its exact axiom audit is recorded separately.

Predictive sufficiency is strictly weaker than autonomy. Let all outputs
and unit costs be identical, with $p\to u$, $q\to v$, $u\to u$, $v\to v$.
Every state has the same behavior, but the representation identifying $p,q$
and distinguishing $u,v$ cannot update autonomously.

For bounded observations the correct state is $(x,b)$ with residual budget
$b$. An affordable transition of cost $c$ leads to $(x',b-c)$; unaffordable
and unavailable requests have the same rejection in this version of the
contract. Applying the theorem to these augmented states gives the lawful
budget quotient. Applying budget-$b$ equivalence repeatedly to bare states
would replenish resources. For example, $p\to u$, $q\to v$, $u\to u$,
$v\to z$, $z\to z$, with unit costs and output one only at $z$, satisfies
$p\simeq_1q$ but $u\not\simeq_1v$.

Effective existence is another boundary. In a computable unary system let
$(M,n)\mapsto(M,n+1)$ and output one exactly when machine $M$ halts within
$n$ steps. Comparison with a forever-zero loop decides nonhalting if
behavioral equivalence is decidable. Thus a computable transition law does
not supply a uniform algorithm for the infinite quotient. This fact has no
implication about the decidability of a particular number-theoretic conjecture.

\begin{theorem}[Exact finite resource threshold]
Suppose $X,A$ are finite, costs are nonnegative rationals, and the budget
interface rejects unavailable or unaffordable requests identically at zero cost,
keeps residual budget hidden, and
reports the exact cost of each successful action. On unordered state pairs
make a terminal edge of weight zero for differing reports; of weight $c$
when only one side enables an action of cost $c$; and of weight
$\min(c,d)$ when both enable it with unequal costs. Equal-cost actions give
edges of that common cost to the successor pair. If $D(x,y)$ is the shortest
distance to a terminal, with $D=+\infty$ if no terminal is reachable, then,
for every rational budget $B\ge0$,
\[
x\not\simeq_B y\quad\Longleftrightarrow\quad D(x,y)\le B.
\]
A finite threshold is attained and has a word witness of at most
$\binom{|X|}{2}$ actions. It can be found without expanding a residual-budget
lattice. Priority for this specialized formulation is unverified.
\end{theorem}
\begin{proof}
Before the first trace discrepancy, both runs execute the same actions with
equal costs, hence follow synchronized pair edges. A discrepancy is either
an output difference, a successful action versus rejection, or unequal
successful cost receipts. Its minimum necessary extra budget is exactly the
corresponding terminal weight. Conversely a terminal path of weight $d$
gives a differing trace at every budget at least $d$: a larger budget can
turn a cost-based affordability difference into unequal cost receipts, but
cannot hide it. Nonnegative weights permit deleting loops. A successful
path never reaches a diagonal pair, which has no possible discrepancy.
Thus a simple shortest terminal path visits at most the stated number of
nondiagonal pairs. The finite rational minimum is attained.
\end{proof}

The observation interface is load-bearing. If one compares only words lawful
at both states, agreement need not be transitive: let $x,y,z$ initially report
zero, let $x$ have a zero-reporting $a$-loop, let $y$ have no $a$-edge, and let
$z$ have an $a$-edge to a one-reporting state. Then $x$ agrees with $y$ and $y$
with $z$ on their common lawful words, while $x,z$ disagree on $a$. Such a
relation cannot be used as a quotient equivalence without an additional
testing contract.

\section{Report gluing versus one lawful history}
Fix a finite private state set $X$, a fixed original-source fiber $X_\sigma$,
initial states $I_\sigma$, a report $r:X_\sigma\to Y$, and relations $R_a$
preserving the original source. Define $y\bar R_a z$ by existence of an
$R_a$ edge between the corresponding private fibers. A projected word is
$w=(y_0,a_0,\ldots,a_{n-1},y_n)$ starting in $r(I_\sigma)$ and using these
projected edges. Its exact private supports are
\[
B_0=I_\sigma\cap r^{-1}(y_0),\qquad
B_{i+1}=R_{a_i}[B_i]\cap r^{-1}(y_{i+1}).
\]
\begin{theorem}[Finite private lifting and obstruction]
A projected word has a private lift if and only if $B_n\ne\varnothing$,
if and only if its complete private layered graph admits a nonnegative unit
flow, if and only if it admits a $0/1$ path flow. Nonexistence is certified by
sets $C_i$ containing $B_0$ initially, closed under the specified successive
transitions, and empty at the final layer.
Every projected word lifts if and only if the nonempty support configurations
reachable from the initial configurations form a closed invariant. If this
fails, a shortest failed word has at most
\[
N_\sigma=\sum_{y\in r(X_\sigma)}(2^{|r^{-1}(y)|}-1)
\]
edges.
\end{theorem}
\begin{proof}
Induction identifies $B_i$ with endpoints of all private lifts of the prefix.
A path gives a unit flow; conservation lets one follow positive-flow edges
from the supersource to the supersink. Time layers prevent cycles. Repeatedly
subtracting the minimum flow on such a path gives a finite convex combination
of path flows, hence integrality for this network contract. Any closed $C_i$
contains $B_i$ inductively, and $C_i=B_i$ provides a cut when $B_n$ is empty.
For incidence columns $e_{\rm head}-e_{\rm tail}$, the indicator of the
source-side closed cut has $p^TA\ge0$ and $p^Tb=-1$, an exact Farkas witness.
There are $N_\sigma$ possible nonempty configurations $(y,B)$. A shortest
failed path cannot repeat one before its final empty update: removing the
loop preserves both report compatibility and the deterministic support
update. The number of nonempty prefixes, and hence of edges, is at most
$N_\sigma$. Exhaustive breadth-first search proves the final alternative.
\end{proof}

The mechanism is classical subset construction, path decomposition and
abstraction checking. The private interfaces are essential. Edges $p\to u$
and $v\to q$ with $r(u)=r(v)$ glue as reports while $u\ne v$ prevents a lift.
Extra marginal constraints can also destroy network integrality: pairwise
binary anticorrelations around an odd cycle admit compatible fractional local
laws but no global deterministic assignment.

Even uniform existential lifting is not an online resolver. With one report,
initial state $p$, edges $p\xrightarrow{a,b}\{u,v\}$,
$u\xrightarrow a\{u,v\}$ and $v\xrightarrow b\{u,v\}$, every finite word
has a lift. Yet any resolver's first choice $u$ or $v$ is defeated by the
next request $b$ or $a$. This separates $\forall w\,\exists\text{lift}$
from one causal strategy for every word, not merely from source authenticity.

\section{Arbitrarily late arithmetic camouflage}
Generalized primes below, in Beurling's framework \cite{beurling}, are positive real generators of a free commutative
monoid; counting retains word multiplicities.
\begin{theorem}[A uniform source nonidentification family]
For every finite $H\ge1$, there is a collision-free rational generalized-prime
source agreeing with ordinary primitive, prime-power and integer data through
$H$, with first real cell length one, $N_B(X)=X+O(\log^2X)$, a prime count
differing from the ordinary one by at most two, and
\[
\zeta_B(s)=\zeta(s)F(s),\qquad \Re s>0,
\]
as a meromorphic continuation from the Euler product on $\Re s>1$,
where $F$ is holomorphic and nowhere zero and $F(1)=1$. Nevertheless its rank
successor and unchanged multiplication by two fail $V_2S=S^2V_2$.
\end{theorem}
\begin{proof}
Choose ordinary primes $p>\max(H,2)$ and $q>p^2$. Put
\[
A=p^2-p+1,\qquad a=\frac{p^2}{p-1},\qquad
b=\frac{qA}{q+p(p-1)},
\]
and replace $p,q$ by $a,b$. Word values are $ma^jb^k$ with $(m,pq)=1$.
Their $p,q$ valuations are $2j,k$, so distinct words have distinct values.
Both replacements exceed $p+1$; consequently the source through $p+1$ is
$1,\ldots,p-1,p+1$. For $x=(p-1)/2$ the difference
\[
\operatorname{rank}(V_2Sx)-\operatorname{rank}(V_2x)-2=-1
\]
exhibits the failed braid, while every bounded operation staying through $H$
is unchanged.

Let $C=(1-1/p)(1-1/q)$ and
$R(y)=\lfloor y\rfloor-\lfloor y/p\rfloor-\lfloor y/q\rfloor+\lfloor y/(pq)\rfloor$.
Then $|R(y)-Cy|\le2$ and direct algebra gives
$(1-1/a)(1-1/b)=C$. Thus
\[
N_B(X)=\sum_{a^jb^k\le X}R(Xa^{-j}b^{-k}).
\]
Writing $J=\max\{j:a^j\le X\}$ and $K=\max\{k:b^k\le X\}$, sum the
rounding error over at most $(J+1)(K+1)$ terms. The omitted geometric tail,
split at $j=J$, is at most $(J+1)(1-1/a)+1$ after multiplication by $CX$.
Hence
\[
|N_B(X)-X|\le2(J+1)(K+1)+(J+1)(1-1/a)+1.
\]
Only two primitive values are replaced, proving the prime-count bound.
On $\Re s>1$ the absolutely convergent Euler products give
\[
F(s)=\frac{(1-p^{-s})(1-q^{-s})}{(1-a^{-s})(1-b^{-s})}.
\]
Every factor is holomorphic and nonzero on $\Re s>0$; ordinary zeta
continuation and $F(1)=1$ prove the rest. In particular the critical-strip
zero divisor and pole residue agree without using any zero ordinate.
\end{proof}

A terminating adaptive protocol using only bounded source-prefix queries has
the same transcript on a sufficiently late member of this family as on
ordinary arithmetic. Even the displayed global summaries do not break that
ambiguity. This does not exclude a finite proof of an authenticated universal
source law. Nor does the construction preserve the standard completed
functional equation, canonical Gamma term or full Weil distribution.
The generalized-prime framework is classical; scholarly priority for this
specific simultaneous construction remains unverified.

\section{Retaining the authentic Gamma tail}
On $I=(-1/2,1/2)$ set $d_j=4j+1$, $\alpha_j=d_j/2$ and
\[
(E_\alpha h)(x)=\int_Ie^{-\alpha|x-y|}h(y)\dd y,
\qquad \Gamma_jh=\frac4{d_j}h-E_{\alpha_j}h.
\]
These are the rungs of the one real-Gamma source in Suzuki's completed form
\cite{suzuki}. They do not replace its signed prime, pole or contact terms.
Let $S_r(J)=\sum_{j\ge J}d_j^{-r}$ and $a_J=16S_3(J)$.

\begin{theorem}[Endpoint-paid tail expansion]
If $J\ge1$ and $h\in H^4(I)\cap H^2_0(I)$, then
\[
\sum_{j\ge J}\Gamma_jh=-a_Jh''+e_J,
\]
where, writing $d=4J+1$ and
$L_J(u,v)=(|u|^2+|v|^2+2|u\bar v|d e^{-d/2})^{1/2}$,
\begin{align*}
\|e_J\|_2\le{}&8L_J(h''(-1/2),h''(1/2))\sqrt{S_3S_4}\\
&+16L_J(-h'''(-1/2),h'''(1/2))\sqrt{S_4S_5}
+64S_5\|h''''\|_2.
\end{align*}
For real endpoint values the cross term may be sharpened from $|u\bar v|$
to $\max(uv,0)$. The quadratic form also satisfies
\[
\langle\Gamma_{\ge J}h,h\rangle
\ge16S_3\|h'\|_2^2-64S_5\|h''\|_2^2.
\]
\end{theorem}
\begin{proof}
Since $h=h'=0$ at the endpoints, two integrations by parts give
$\Gamma_jh=-\alpha_j^{-2}E_{\alpha_j}h''$. Two more yield
\begin{align*}
\Gamma_jh={}&-16d_j^{-3}h''-16d_j^{-4}E_{\alpha_j}h''''\\
&+e^{-\alpha_j(x+1/2)}(8d_j^{-3}h''(-1/2)-16d_j^{-4}h'''(-1/2))\\
&+e^{-\alpha_j(1/2-x)}(8d_j^{-3}h''(1/2)+16d_j^{-4}h'''(1/2)).
\end{align*}
The exponential operator has norm at most $4/d_j$. The squared norm of an
endpoint pair is
$(|u|^2+|v|^2)(1-e^{-d_j})/d_j+2\Re(u\bar v)e^{-d_j/2}$.
Since $d_j e^{-d_j/2}$ decreases for $d_j\ge5$, the stated common pair bound
follows. Sum the rung norms and use
$\sum d_j^{-7/2}\le\sqrt{S_3S_4}$ and
$\sum d_j^{-9/2}\le\sqrt{S_4S_5}$.
For the quadratic bound, zero extension belongs to $H^2(\RR)$.
The Fourier multiplier of a rung is
$16\xi^2/[d_j(d_j^2+4\xi^2)]$; apply $(1+t)^{-1}\ge1-t$ and Plancherel.
\end{proof}

The endpoint contribution is $O(J^{-5/2})$ and the bulk remainder
$O(J^{-4})$, for fixed $h$. The estimate does not impose the false extra
conditions $h''=h'''=0$ at the endpoints. Exact rational series enclosures
and outward square roots implement it. It reduced tail-action uncertainty
on all four tested $J=24$ columns,
but the retained action also moves the residual center: a smaller error
bound alone need not improve a signed Schur lower bound.

Let $F_c$ be the original four-dimensional low vector with coefficients $c$,
$E$ the original primitive high subspace, $Q_E$ its orthogonal projection,
and $H_\infty$ the original self-adjoint signed source operator. For a trial
$h_c=F_c-x_c$ with $x_c\in E$ in the valid common form domains and $h_c$ in
the operator domain, and inherited original high coercivity $D\ge\kappa I$
in the correct primitive Hilbert space, the valid bound is
\[
S_4[c]\ge q(h_c)-\kappa^{-1}\|Q_EH_\infty h_c\|_2^2.
\]
The new estimate controls part of this expression; it proves neither its sign
nor every-radius positivity. Paired source trials remain ambiguous. The
complete signed source, one real Gamma, both polar signs, contact and genuine
prime half-density remain fixed throughout.

\section{Reproducibility and status}
This paper accompanies the branch \texttt{research/academic-transmutation-2026-10-18},
based on draft PR177 commit \texttt{f309d9104ba480c1999c962af1df7c50f19e1bfe}.
The original signed source blob is
\texttt{0cc8510b71d178fc79e8d410a86d8b6014cf9178}.
The accompanying dossiers supply complete domains, exact scripts, independent
finite checkers, mutation controls, primary-source comparisons and run receipts.
Lean 4.14.0 checks the quotient/path fragment, not the analytic or arithmetic
theorems, not the Python programs and not the whole framework. The finite
certificates establish statements about their supplied tables. The source
nonidentification and Gamma estimates are self-contained paper derivations.
All academic priority claims remain unverified. RH remains open.

\begin{thebibliography}{9}
\bibitem{rutten} J. J. M. M. Rutten, Universal coalgebra: a theory of systems,
\emph{Theoretical Computer Science} 249 (2000), 3--80.
\bibitem{silva} A. Silva, F. Bonchi, M. Bonsangue and J. Rutten,
Generalizing determinization from automata to coalgebras,
\url{https://arxiv.org/abs/1302.1046}.
\bibitem{beurling} A. Beurling, Analyse de la loi asymptotique de la distribution
des nombres premiers generalises. I, \emph{Acta Mathematica} 68 (1937), 255--291.
\bibitem{suzuki} M. Suzuki, Weil's quadratic form via the screw function,
version 3, \url{https://arxiv.org/abs/2606.09096} (2026).
\end{thebibliography}
\end{document}
